% Options for packages loaded elsewhere
\PassOptionsToPackage{unicode}{hyperref}
\PassOptionsToPackage{hyphens}{url}
%
\documentclass[
  ignorenonframetext,
]{beamer}
\usepackage{pgfpages}
\setbeamertemplate{caption}[numbered]
\setbeamertemplate{caption label separator}{: }
\setbeamercolor{caption name}{fg=normal text.fg}
\beamertemplatenavigationsymbolsempty
% Prevent slide breaks in the middle of a paragraph
\widowpenalties 1 10000
\raggedbottom
\setbeamertemplate{part page}{
  \centering
  \begin{beamercolorbox}[sep=16pt,center]{part title}
    \usebeamerfont{part title}\insertpart\par
  \end{beamercolorbox}
}
\setbeamertemplate{section page}{
  \centering
  \begin{beamercolorbox}[sep=12pt,center]{section title}
    \usebeamerfont{section title}\insertsection\par
  \end{beamercolorbox}
}
\setbeamertemplate{subsection page}{
  \centering
  \begin{beamercolorbox}[sep=8pt,center]{subsection title}
    \usebeamerfont{subsection title}\insertsubsection\par
  \end{beamercolorbox}
}
\AtBeginPart{
  \frame{\partpage}
}
\AtBeginSection{
  \ifbibliography
  \else
    \frame{\sectionpage}
  \fi
}
\AtBeginSubsection{
  \frame{\subsectionpage}
}

\usepackage{amsmath,amssymb}
\usepackage{iftex}
\ifPDFTeX
  \usepackage[T1]{fontenc}
  \usepackage[utf8]{inputenc}
  \usepackage{textcomp} % provide euro and other symbols
\else % if luatex or xetex
  \usepackage{unicode-math}
  \defaultfontfeatures{Scale=MatchLowercase}
  \defaultfontfeatures[\rmfamily]{Ligatures=TeX,Scale=1}
\fi
\usepackage{lmodern}
\ifPDFTeX\else  
    % xetex/luatex font selection
\fi
% Use upquote if available, for straight quotes in verbatim environments
\IfFileExists{upquote.sty}{\usepackage{upquote}}{}
\IfFileExists{microtype.sty}{% use microtype if available
  \usepackage[]{microtype}
  \UseMicrotypeSet[protrusion]{basicmath} % disable protrusion for tt fonts
}{}
\makeatletter
\@ifundefined{KOMAClassName}{% if non-KOMA class
  \IfFileExists{parskip.sty}{%
    \usepackage{parskip}
  }{% else
    \setlength{\parindent}{0pt}
    \setlength{\parskip}{6pt plus 2pt minus 1pt}}
}{% if KOMA class
  \KOMAoptions{parskip=half}}
\makeatother
\usepackage{xcolor}
\newif\ifbibliography
\setlength{\emergencystretch}{3em} % prevent overfull lines
\setcounter{secnumdepth}{-\maxdimen} % remove section numbering

\usepackage{color}
\usepackage{fancyvrb}
\newcommand{\VerbBar}{|}
\newcommand{\VERB}{\Verb[commandchars=\\\{\}]}
\DefineVerbatimEnvironment{Highlighting}{Verbatim}{commandchars=\\\{\}}
% Add ',fontsize=\small' for more characters per line
\usepackage{framed}
\definecolor{shadecolor}{RGB}{241,243,245}
\newenvironment{Shaded}{\begin{snugshade}}{\end{snugshade}}
\newcommand{\AlertTok}[1]{\textcolor[rgb]{0.68,0.00,0.00}{#1}}
\newcommand{\AnnotationTok}[1]{\textcolor[rgb]{0.37,0.37,0.37}{#1}}
\newcommand{\AttributeTok}[1]{\textcolor[rgb]{0.40,0.45,0.13}{#1}}
\newcommand{\BaseNTok}[1]{\textcolor[rgb]{0.68,0.00,0.00}{#1}}
\newcommand{\BuiltInTok}[1]{\textcolor[rgb]{0.00,0.23,0.31}{#1}}
\newcommand{\CharTok}[1]{\textcolor[rgb]{0.13,0.47,0.30}{#1}}
\newcommand{\CommentTok}[1]{\textcolor[rgb]{0.37,0.37,0.37}{#1}}
\newcommand{\CommentVarTok}[1]{\textcolor[rgb]{0.37,0.37,0.37}{\textit{#1}}}
\newcommand{\ConstantTok}[1]{\textcolor[rgb]{0.56,0.35,0.01}{#1}}
\newcommand{\ControlFlowTok}[1]{\textcolor[rgb]{0.00,0.23,0.31}{\textbf{#1}}}
\newcommand{\DataTypeTok}[1]{\textcolor[rgb]{0.68,0.00,0.00}{#1}}
\newcommand{\DecValTok}[1]{\textcolor[rgb]{0.68,0.00,0.00}{#1}}
\newcommand{\DocumentationTok}[1]{\textcolor[rgb]{0.37,0.37,0.37}{\textit{#1}}}
\newcommand{\ErrorTok}[1]{\textcolor[rgb]{0.68,0.00,0.00}{#1}}
\newcommand{\ExtensionTok}[1]{\textcolor[rgb]{0.00,0.23,0.31}{#1}}
\newcommand{\FloatTok}[1]{\textcolor[rgb]{0.68,0.00,0.00}{#1}}
\newcommand{\FunctionTok}[1]{\textcolor[rgb]{0.28,0.35,0.67}{#1}}
\newcommand{\ImportTok}[1]{\textcolor[rgb]{0.00,0.46,0.62}{#1}}
\newcommand{\InformationTok}[1]{\textcolor[rgb]{0.37,0.37,0.37}{#1}}
\newcommand{\KeywordTok}[1]{\textcolor[rgb]{0.00,0.23,0.31}{\textbf{#1}}}
\newcommand{\NormalTok}[1]{\textcolor[rgb]{0.00,0.23,0.31}{#1}}
\newcommand{\OperatorTok}[1]{\textcolor[rgb]{0.37,0.37,0.37}{#1}}
\newcommand{\OtherTok}[1]{\textcolor[rgb]{0.00,0.23,0.31}{#1}}
\newcommand{\PreprocessorTok}[1]{\textcolor[rgb]{0.68,0.00,0.00}{#1}}
\newcommand{\RegionMarkerTok}[1]{\textcolor[rgb]{0.00,0.23,0.31}{#1}}
\newcommand{\SpecialCharTok}[1]{\textcolor[rgb]{0.37,0.37,0.37}{#1}}
\newcommand{\SpecialStringTok}[1]{\textcolor[rgb]{0.13,0.47,0.30}{#1}}
\newcommand{\StringTok}[1]{\textcolor[rgb]{0.13,0.47,0.30}{#1}}
\newcommand{\VariableTok}[1]{\textcolor[rgb]{0.07,0.07,0.07}{#1}}
\newcommand{\VerbatimStringTok}[1]{\textcolor[rgb]{0.13,0.47,0.30}{#1}}
\newcommand{\WarningTok}[1]{\textcolor[rgb]{0.37,0.37,0.37}{\textit{#1}}}

\providecommand{\tightlist}{%
  \setlength{\itemsep}{0pt}\setlength{\parskip}{0pt}}\usepackage{longtable,booktabs,array}
\usepackage{calc} % for calculating minipage widths
\usepackage{caption}
% Make caption package work with longtable
\makeatletter
\def\fnum@table{\tablename~\thetable}
\makeatother
\usepackage{graphicx}
\makeatletter
\newsavebox\pandoc@box
\newcommand*\pandocbounded[1]{% scales image to fit in text height/width
  \sbox\pandoc@box{#1}%
  \Gscale@div\@tempa{\textheight}{\dimexpr\ht\pandoc@box+\dp\pandoc@box\relax}%
  \Gscale@div\@tempb{\linewidth}{\wd\pandoc@box}%
  \ifdim\@tempb\p@<\@tempa\p@\let\@tempa\@tempb\fi% select the smaller of both
  \ifdim\@tempa\p@<\p@\scalebox{\@tempa}{\usebox\pandoc@box}%
  \else\usebox{\pandoc@box}%
  \fi%
}
% Set default figure placement to htbp
\def\fps@figure{htbp}
\makeatother

\makeatletter
\@ifpackageloaded{caption}{}{\usepackage{caption}}
\AtBeginDocument{%
\ifdefined\contentsname
  \renewcommand*\contentsname{Table of contents}
\else
  \newcommand\contentsname{Table of contents}
\fi
\ifdefined\listfigurename
  \renewcommand*\listfigurename{List of Figures}
\else
  \newcommand\listfigurename{List of Figures}
\fi
\ifdefined\listtablename
  \renewcommand*\listtablename{List of Tables}
\else
  \newcommand\listtablename{List of Tables}
\fi
\ifdefined\figurename
  \renewcommand*\figurename{Figure}
\else
  \newcommand\figurename{Figure}
\fi
\ifdefined\tablename
  \renewcommand*\tablename{Table}
\else
  \newcommand\tablename{Table}
\fi
}
\@ifpackageloaded{float}{}{\usepackage{float}}
\floatstyle{ruled}
\@ifundefined{c@chapter}{\newfloat{codelisting}{h}{lop}}{\newfloat{codelisting}{h}{lop}[chapter]}
\floatname{codelisting}{Listing}
\newcommand*\listoflistings{\listof{codelisting}{List of Listings}}
\makeatother
\makeatletter
\makeatother
\makeatletter
\@ifpackageloaded{caption}{}{\usepackage{caption}}
\@ifpackageloaded{subcaption}{}{\usepackage{subcaption}}
\makeatother

\usepackage{bookmark}

\IfFileExists{xurl.sty}{\usepackage{xurl}}{} % add URL line breaks if available
\urlstyle{same} % disable monospaced font for URLs
\hypersetup{
  pdftitle={PSY 3307 Slides},
  pdfauthor={Jonathan A. Pedroza, PhD},
  hidelinks,
  pdfcreator={LaTeX via pandoc}}


\title{PSY 3307 Slides}
\author{Jonathan A. Pedroza, PhD}
\date{}

\begin{document}
\frame{\titlepage}

\renewcommand*\contentsname{Table of contents}
\begin{frame}[allowframebreaks]
  \frametitle{Table of contents}
  \setcounter{tocdepth}{3}
  \tableofcontents
\end{frame}

\begin{frame}{Week 1 - Introduction \& Review}
\phantomsection\label{week-1---introduction-review}
\end{frame}

\begin{frame}{Agenda}
\phantomsection\label{agenda}
\begin{itemize}
\tightlist
\item
  Introductions
\item
  Syllabus
\item
  Review
\item
  Types of Research
\item
  Steps of Conducting Statistics Within the Research Process
\item
  Levels of Measurement

  \begin{itemize}
  \tightlist
  \item
    Measurement Error
  \end{itemize}
\item
  Measurement Scales
\end{itemize}
\end{frame}

\begin{frame}{About Me}
\phantomsection\label{about-me}
\begin{itemize}
\tightlist
\item
  My name is Jonathan A. Pedroza (he/him)

  \begin{itemize}
  \tightlist
  \item
    please refer to me as \textbf{JP} or Jon
  \end{itemize}
\item
  First-generation Graduate

  \begin{itemize}
  \tightlist
  \item
    CPP Alumni (2014)
  \item
    CSUDH (2017)
  \item
    University of Oregon (2021)
  \end{itemize}
\end{itemize}
\end{frame}

\begin{frame}{About Me}
\phantomsection\label{about-me-1}
\begin{figure}[H]

{\centering \pandocbounded{\includegraphics[keepaspectratio]{../additional_resources/hike.jpg}}

}

\caption{uo grad}

\end{figure}%
\end{frame}

\begin{frame}{About Me}
\phantomsection\label{about-me-2}
\begin{figure}[H]

{\centering \pandocbounded{\includegraphics[keepaspectratio]{../additional_resources/hike.jpg}}

}

\caption{hiking}

\end{figure}%
\end{frame}

\begin{frame}{About Me}
\phantomsection\label{about-me-3}
\begin{figure}[H]

{\centering \pandocbounded{\includegraphics[keepaspectratio]{../additional_resources/lucy.jpg}}

}

\caption{lucy}

\end{figure}%
\end{frame}

\begin{frame}[fragile]{Nerdy Stuff}
\phantomsection\label{nerdy-stuff}
\begin{itemize}
\tightlist
\item
  Passion/Nerdy Topics:

  \begin{itemize}
  \tightlist
  \item
    I love creating data visualizations in \texttt{R}
  \item
    health inequities in communities of color
  \item
    examining environmental variables that influence health behaviors
  \end{itemize}
\end{itemize}
\end{frame}

\begin{frame}{Nerdy Stuff}
\phantomsection\label{nerdy-stuff-1}
\begin{figure}[H]

{\centering \pandocbounded{\includegraphics[keepaspectratio]{../additional_resources/non_passing_rate_visualization.png}}

}

\caption{non\_passing\_rate}

\end{figure}%
\end{frame}

\begin{frame}{Break-out Rooms - Tell Me About Yourself}
\phantomsection\label{break-out-rooms---tell-me-about-yourself}
\begin{itemize}
\item
  Name \& pronouns
\item
  What do you want to do after you graduate?
\item
  What is your experience with statistics?
\item
  What is your experience with SPSS?
\item
  Something that you are passionate/nerdy about
\end{itemize}
\end{frame}

\begin{frame}{Learning Objectives}
\phantomsection\label{learning-objectives}
\begin{itemize}
\item
  Use SPSS to compute statistics
\item
  Read/Interpret Descriptive \& Inferential Statistics
\item
  Understand \& Solve Formulas for Statistical Tests
\item
  Understand Sampling, Statistical Null Hypothesis Testing, Effect
  Sizes, Confidence Intervals
\item
  Demonstrate ability to apply these skills to test hypotheses using
  real data
\end{itemize}
\end{frame}

\begin{frame}{Topics Covered}
\phantomsection\label{topics-covered}
\begin{itemize}
\tightlist
\item
  Descriptive Statistics
\item
  Probability
\item
  Random Selection
\item
  Normal Curve \& Sampling Distribution
\item
  Hypothesis Testing
\item
  Rules for Statistical Decisions (Types I and II Errors)
\item
  Sample Size Importance (One- vs Two-tailed Tests)
\item
  Z-scores, z-ratio, and t-ratio (single group, independent-samples
  t-test, paired)
\item
  Analysis of Variance (One-way, Two-way, Repeated-measure, mixed)
\item
  Confidence Intervals/Error bars in plots
\item
  Most common plots (histograms, boxplots, scatterplots)
\item
  Correlation \& Regression Analysis
\item
  Chi-square (One-way, Two-way)
\item
  Non-parametric statistics
\end{itemize}
\end{frame}

\begin{frame}{You Might be Feeling Like\ldots{}}
\phantomsection\label{you-might-be-feeling-like}
\begin{figure}[H]

{\centering \pandocbounded{\includegraphics[keepaspectratio]{Slides_files/mediabag/giphy.gif}}

}

\caption{spongebob running}

\end{figure}%
\end{frame}

\begin{frame}{What I Need From You}
\phantomsection\label{what-i-need-from-you}
\begin{itemize}
\item
  Cameras are optional

  \begin{itemize}
  \tightlist
  \item
    It helps me see if you are confused
  \end{itemize}
\item
  Unmute yourself and interrupt me
\item
  put in the chat that you are confused
\item
  private message me in the chat that you are confused
\end{itemize}
\end{frame}

\begin{frame}{Week 1 - Research Methods}
\phantomsection\label{week-1---research-methods}
\end{frame}

\begin{frame}{Designs}
\phantomsection\label{designs}
\begin{itemize}
\item
  \textbf{Design} is how a study is organized
\item
  Experimental (manipulate your independent variable to see what effect
  it has on the dependent variable)

  \begin{itemize}
  \tightlist
  \item
    True Experiment
  \item
    Quasi Experiment
  \item
    Natural Experiment
  \end{itemize}
\item
  Correlational (don't manipulate anything; just look at the
  relationship between two variables)

  \begin{itemize}
  \tightlist
  \item
    Correlation
  \item
    Regression
  \end{itemize}
\end{itemize}
\end{frame}

\begin{frame}{Differences in Population \& Sample}
\phantomsection\label{differences-in-population-sample}
\begin{itemize}
\item
  \textbf{Population} is a large group of individuals, which a law of
  nature applies
\item
  \textbf{Sample} is a group of a given population intended to represent
  that population
\item
  \textbf{Participants} are those measured in a sample. Participants
  always refers to individuals (e.g., students, children, prisoners)
\item
  A sample is \emph{supposed to generalize} to a given population
\item
  We never use subjects in the social sciences anymore, they are
  participants
\end{itemize}
\end{frame}

\begin{frame}{Statistics vs Parameters}
\phantomsection\label{statistics-vs-parameters}
\begin{itemize}
\item
  \textbf{Statistics} use English letters to get values for a sample
\item
  \textbf{Parameter} use greek letters for values of a population
\end{itemize}

Statistic = Sample

Parameter = Population
\end{frame}

\begin{frame}{Straight into Terms}
\phantomsection\label{straight-into-terms}
\begin{itemize}
\item
  \textbf{raw score} is the score given to a participant
\item
  \textbf{Frequency} denoted as f; number of times a score occurs/is
  counted
\end{itemize}

\emph{Note.} Not F, that is something completely different.

\begin{itemize}
\tightlist
\item
  \textbf{Frequency Distribution} is a distribution of each score and
  the number of times the score has occurred/is counted
\end{itemize}
\end{frame}

\begin{frame}{Visualizing Frequencies}
\phantomsection\label{visualizing-frequencies}
\begin{itemize}
\item
  Best way of visualizing frequencies is by using a bar graph
\item
  \textbf{Bar graph} graph with vertical bar over each nominal/ordinal
  category
\item
  \textbf{Side Note} A pie graph will always be inferior to a bar
  graph/any other visual
\item
  \textbf{Histogram} is a frequency graph used for interval or ratio
  scores
\item
  \textbf{Frequency Polygon} similar to a histogram, which shows data
  points connected with straight lines
\item
  \textbf{Grouped Distributions} put continuous data into categories

  \begin{itemize}
  \tightlist
  \item
    The next slide has data ranging from 6-10 in coffee uniformity; I
    could lump them as anything below a perfect 10 (6-9) and then 10 as
    a different category
  \end{itemize}
\end{itemize}
\end{frame}

\begin{frame}{Normal Distribution}
\phantomsection\label{normal-distribution}
\begin{figure}[H]

{\centering \pandocbounded{\includegraphics[keepaspectratio]{Slides_files/mediabag/graphic-1.gif}}

}

\caption{paranormal distribution}

\end{figure}%
\end{frame}

\begin{frame}{Why Is It a Normal Distribution?}
\phantomsection\label{why-is-it-a-normal-distribution}
\begin{itemize}
\item
  \textbf{Normal curve} is often called the bell-shaped curve; is
  symmetrical
\item
  \textbf{Normal Distribution} same thing as normal curve; represents
  the population because if you have enough data you will get a normal
  distribution (central limit theorem); if your data looks like this in
  a histogram, you're in good shape
\item
  \textbf{Distribution Tail} has two tails; these will be more important
  for statistics
\end{itemize}
\end{frame}

\begin{frame}{Skewed Distributions}
\phantomsection\label{skewed-distributions}
\begin{itemize}
\item
  \textbf{Negative Skew} is like a finger pointing left; not normal and
  is asymmetrical; indicates higher frequency of middle and higher
  scores; no low frequency in higher scores
\item
  \textbf{Positive Skew} is like a finger pointing right; not normal and
  is asymmetrical; indicates higher frequency of low and middle scores;
  no low frequency in lower scores
\item
  Some thresholds are that if you have a skewness value of +-2 or +-3
  then you're good to use that variable like it is.
\item
  \textbf{Kurtosis} is when your frequency are really skinny and tall or
  really flat and wide
\end{itemize}
\end{frame}

\begin{frame}{Bimodal Distribution}
\phantomsection\label{bimodal-distribution}
\begin{itemize}
\tightlist
\item
  \textbf{Bimodal Distribution} is when your distribution has two humps
  with a valley in the middle; high frequencies both below and above the
  middle of the plot
\end{itemize}
\end{frame}

\begin{frame}{Some Notes From JP}
\phantomsection\label{some-notes-from-jp}
Deciding what is positively and negatively skewed from visuals alone is
not good enough

You'll want to look into your descriptive statistics and look at
skewness and kurtosis to make sure they are within +-3

Some statistics can handle some skewness and kurtosis

Other times you'll have to transform the variable using fancy methods
that we will not talk about in this class.
\end{frame}

\begin{frame}[fragile]{Relative Frequency}
\phantomsection\label{relative-frequency}
\begin{itemize}
\tightlist
\item
  \textbf{Relative Frequency} the proportion of times a score occures/is
  counted in the distribution
\end{itemize}

\[ Relative\;frequency = f/N\]

Here, f is the frequency of one category of a nominal variable. N is the
total number of observations for all the categories of that variable.

\begin{Shaded}
\begin{Highlighting}[]
\NormalTok{f }\OtherTok{=} \DecValTok{37}
\NormalTok{N }\OtherTok{=} \DecValTok{2000}

\NormalTok{relative\_frequency }\OtherTok{=}\NormalTok{ f}\SpecialCharTok{/}\NormalTok{N}
\NormalTok{relative\_frequency}
\end{Highlighting}
\end{Shaded}

\begin{verbatim}
[1] 0.0185
\end{verbatim}

\begin{Shaded}
\begin{Highlighting}[]
\NormalTok{percent }\OtherTok{=}\NormalTok{ relative\_frequency}\SpecialCharTok{*}\DecValTok{100}
\NormalTok{percent}
\end{Highlighting}
\end{Shaded}

\begin{verbatim}
[1] 1.85
\end{verbatim}
\end{frame}

\begin{frame}{Relative Frequency Using Normal Curve}
\phantomsection\label{relative-frequency-using-normal-curve}
\begin{itemize}
\item
  \textbf{Proportion of Area under the Curve} is the proportion of total
  area under the normal curve
\item
  \textbf{Percentile} is the percentage of all scores in the sample
  below a particular score
\item
  \textbf{Cumulative Frequency} is the number of scores in the data at
  or below a particular score
\end{itemize}
\end{frame}

\begin{frame}{Greek Terms}
\phantomsection\label{greek-terms}
\[\Sigma = Sum\;of\;scores\]

\[\Sigma\;X = Sum\;Of\;X\]
\end{frame}

\begin{frame}{Greek Terms}
\phantomsection\label{greek-terms-1}
\begin{itemize}
\item
  Both of these are stating that we are adding all of the data points
  together.
\item
  \textbf{Sum of X} is the sum of the scores in a sample

  \begin{itemize}
  \tightlist
  \item
    X is just another way to say all the data points for the variable.
  \end{itemize}
\end{itemize}
\end{frame}

\begin{frame}{Central Tendency}
\phantomsection\label{central-tendency}
\begin{itemize}
\item
  Concept that as statisticians every person is just a data point
\item
  We are interested in the central score

  \begin{itemize}
  \tightlist
  \item
    We are interested in how much a person is away from that central
    score
  \item
    But when it comes to statistics we like to group together our
    participants' scores/values
  \end{itemize}
\item
  \textbf{Measures of Central Tendency} are statistics that summarize
  the location of a distribution on a variable by indicating where the
  center is
\item
  A Normal distribution will have the central point right down the
  center
\item
  A skewed distribution will have the central point where the frequency
  of scores is the highest
\end{itemize}
\end{frame}

\begin{frame}{The Mode}
\phantomsection\label{the-mode}
.pull-left{[}
\pandocbounded{
\includegraphics[keepaspectratio]{Slides_files/mediabag/prj6o.jpg}}{]}

.pull-right{[} * Value/score with the highest frequency

\begin{itemize}
\item
  Essentially useless in statistics
\item
  \textbf{Unimodal distributions} distribution with only one hump; one
  value is the mode
\item
  \textbf{Bimodal distributions} distribution with two humps; two values
  with the highest frequency; two modes
\end{itemize}
\end{frame}

\begin{frame}{The Median}
\phantomsection\label{the-median}
.pull-left{[}
\pandocbounded{
\includegraphics[keepaspectratio]{Slides_files/mediabag/Thumbnail_expert_mic.png}}{]}

.pull-right{[} * \textbf{Median} is the middle value/score; the 50th
percentile

\begin{itemize}
\item
  Unlike the mode, it will always be close to the middle of a
  distribution
\item
  You'll only ever have one median
\item
  The symbol is: \[Mdn = Median\] {]}
\end{itemize}
\end{frame}

\begin{frame}{The Median}
\phantomsection\label{the-median-1}
\begin{itemize}
\item
  Preferred for ordinal/ordered data
\item
  Not the best option for normally distributed interval \& ratio scores
\item
  Is more reliable when dealing with skewed data
\item
  \emph{Important Note} If you have an even number of scores/values,
  then you will add the two middle values and divide by 2
\end{itemize}
\end{frame}

\begin{frame}[fragile]{The Mean/Average}
\phantomsection\label{the-meanaverage}
\begin{itemize}
\item
  \textbf{Mean} is the score located at the \texttt{mathematical} center
  of a distribution
\item
  Xbar is often used for the mean

  \begin{itemize}
  \tightlist
  \item
    \[ \overline{X} \]
  \end{itemize}
\item
  \[\overline{X} = \frac{\Sigma\;X}{N} \] is the formula to calculate
  the mean.

  \begin{itemize}
  \tightlist
  \item
    Xbar is the Sum of X/Scores divided by the total number of
    observations/scores/values
  \end{itemize}
\item
  Calculate the mean for interval and ratio scales

  \begin{itemize}
  \tightlist
  \item
    The mean of ordinal/ordered data makes no sense
  \end{itemize}
\item
  Basis for most \texttt{inferential\ statistics}
\end{itemize}
\end{frame}

\begin{frame}{When to Use Each Measure}
\phantomsection\label{when-to-use-each-measure}
\begin{itemize}
\tightlist
\item
  The Median more accurately describes/summarizes skewed data compared
  to the mean

  \begin{itemize}
  \tightlist
  \item
    The mean will be pulled toward the extreme tail of the distribution
  \end{itemize}
\item
  Normal distribution then use the mean as the best measure to
  describe/summarize data
\end{itemize}
\end{frame}

\begin{frame}{The Mean in Research}
\phantomsection\label{the-mean-in-research}
\begin{itemize}
\item
  Most statistics revolve around the mean
\item
  Can't just trust descriptive statistics like the mean
\end{itemize}
\end{frame}

\begin{frame}{Deviation}
\phantomsection\label{deviation}
\begin{itemize}
\item
  The distance a participant's score/value is from the mean/average
\item
  Deviations can be positive or negative

  \begin{itemize}
  \tightlist
  \item
    participants can score lower (negative) than the mean and higher
    (positive) than the mean
  \end{itemize}
\item
  To get the deviation, you subtract the mean from each participant's
  score
\item
  \[ X - \overline{X} \]
\item
  The larger the value the farther away from the mean the score/value is
\item
  \textbf{Sum of the deviations around the mean} is the sum of all
  differences between the scores and the mean
\end{itemize}
\end{frame}

\begin{frame}{Looking to the Future}
\phantomsection\label{looking-to-the-future}
\begin{itemize}
\item
  Deviations is the start for upcoming lectures and statistical tests,
  especially the sum of the deviations
\item
  \[ \Sigma(X - \overline{X}) \]
\item
  If the sum of the deviations is 0 then that means your math is good
\item
  Deviation of each score/value from the mean is often referred to as
  error/residual in statistical tests
\item
  Correlational designs use the mean of IV and the mean of DV to look
  for a relationship between the two variables
\item
  Experiments compare the two or more groups (IV) and the relationship
  with the mean value/score of the DV
\end{itemize}
\end{frame}

\begin{frame}{Visuals}
\phantomsection\label{visuals}
\begin{itemize}
\item
  Line Graphs are good for showing progress or change in time
\item
  To show group differences (nominal or ordinal IV), bar graphs are the
  norm
\item
  Scatterplots are best for continuous IV and continuous DV (interval or
  ratio)
\end{itemize}
\end{frame}

\begin{frame}{Describing the Population Mean}
\phantomsection\label{describing-the-population-mean}
\begin{itemize}
\item
  \[ \mu = Population\;Mean \]
\item
  If you test a population then you would use mu instead of xbar
\item
  \[ \mu = \frac{\Sigma\;X}{N} \]
\end{itemize}
\end{frame}




\end{document}
